\section{Definizione del Global Alignment problem e dei prerequisiti essenziali}

\subsection{Rappresentazione delle sequenze biologiche}

Il DNA è costituito da quattro tipi di nucleotidi, caratterizzati dalle seguenti basi azotate:
\begin{itemize}
    \item Adenina, $A$;
    \item Citosina, $C$;
    \item Guanina, $G$;
    \item Timina, $T$.
\end{itemize}
I due filamenti della molecola sono complementari e antiparalleli, con appaiamenti
\[
\begin{matrix}
    A \longleftrightarrow T & , & C\longleftrightarrow G
\end{matrix}
\]

Dal punto di vista computazionale, un singolo filamento può essere rappresentato come una stringa sull'alfabeto
\[
    \Sigma_{\text{DNA}} = \{A,C,G,T\}.
\]

Quindi una stringa di lunghezza $n$ ha forma
\[
    x=x_1x_2\cdots x_n,\qquad x_i\in\Sigma_{DNA}\qquad \forall i \in\{1,\ldots,n\}.
\]

Una \emph{substring} è una porzione contigua di una stringa. Data una stringa di lunghezza $n$ esiste una substring di lunghezza $k$ per $k\leq n$.

Un \emph{$k$-mero} è una substring di lunghezza $k\geq1$.

Una \emph{subsequence} è invece una nuova stringa ottenuta eliminando 0 o più elementi della stringa originale mantenendone l'ordine relativo.

Mostriamo un esempio: sia $s$ una stringa che rappresenta una sequenza genomica,
\[
    s:= \texttt{ACCT},
\]
allora $\texttt{AC}$ può essere sia una substring che una subsequence, mentre $\texttt{AT}$ è definibile soltanto come subsequence.

Due sequenze genomiche omologhe possono differire tra loro a causa di mutazioni. Consideriamo tre eventi possibili associati a tali mutazioni:

\[
\begin{array}{r c l}
    \texttt{AC}  & \xrightarrow{\text{sostituzione}} & \texttt{AA} \\
    \texttt{AC}  & \xrightarrow{\text{inserzione}}   & \texttt{AAC} \\ 
    \texttt{AAC} & \xrightarrow{\text{delezione}}    & \texttt{AC}
\end{array}
\]

Gli ultimi due eventi verranno ogni tanto trattati come appartenenti a un unico tipo, dal nome \emph{indel}.

Supponiamo, quindi, di avere due sequenze genomiche omologhe che differiscono a causa di mutazioni genetiche; come possiamo quantificare e
individuare la somiglianza tra le due sequenze?

\subsection{Allineamento e modello di scoring}

Stiamo quindi iniziando a definire il problema di nostro interesse.

Prima però chiariamo velocemente l'idea di problema e algoritmo in generale.

Dato che un nostro obiettivo è quello di misurare la differenza tra due sequenze genetiche, trattiamo adesso la \emph{distanza di Hamming}, che è lo strumento più semplice per valutare la somiglianza tra due stringhe.
Siano $x,y$ due stringhe della stessa lunghezza, allora la distanza di Hamming è
\[
    d_H(x,y) := |\{i:x_i\neq y_i\}|,
\]
ovvero il numero di posizioni in cui due stringhe differiscono.
La distanza di Hamming è però applicabile soltanto a stringhe della stessa lunghezza e non modella adeguatamente inserzioni e delezioni, che possono
inoltre alterare la corrispondenza tra le posizioni delle due sequenze.

Abbiamo quindi bisogno di una misura più generale.
Definiamo allora l'\emph{edit distance}, o anche distanza di Levenshtein, $d(v,w)$, come il numero minimo di operazioni elementari
di inserzione, delezione e sostituzione necessarie per trasformare una sequenza $v$ in una sequenza $w$.

L'edit distance fornisce una misura della dissimilarità tra due sequenze. Per rappresentare esplicitamente le corrispondenze tra i loro simboli
introduciamo ora il concetto di \textbf{allineamento}.

Siano $v$ e $w$ due sequenze. Un allineamento di $v$ e $w$ è una matrice 
\[
    \mathcal{A}\in(\Sigma_{\text{DNA}}\cup \{ - \})^{2\times L},
\]
con,
\[
    \max \{ |v|,|w| \}\leq L \leq |v| + |w|,
\]
tale che, eliminando i simboli di gap dalla prima riga, si ottiene $v$, mentre eliminandoli dalla seconda si ottiene $w$.
Non sono ammesse colonne contenenti due gap.

Ogni colonna di un allineamento corrisponde a uno dei seguenti casi:
\begin{itemize}
    \item match;
    \item mismatch (sostituzione);
    \item inserzione;
    \item delezione;
\end{itemize}

Visualizziamo cosa è stato appena detto: siano $v_{aligned}$ e $w_{aligned}$ le sequenze allineate e $\mathcal{A}$ il loro allineamento, allora, \[ \begin{matrix}
        v_{\text{aligned}} & : & A & C & T & - & G & C & A \\
        w_{\text{aligned}} & : & A & - & T & C & G & C & A
    \end{matrix}\enspace = \enspace\mathcal{A}
\]

Per valutare la similarità tra simboli introduciamo una funzione di \textbf{scoring},
\[
    \delta : ((\Sigma_{\text{DNA}}\cup \{ - \})^{2}\setminus \{ (-,-) \}) \to \mathbb{R},
\]
che, dati due simboli $a$ e $b$, restituisce il contributo di score associato al loro allineamento.

Una funzione di scoring (anche detta modello) per una sequenza di DNA è generalmente definita attribuendo un premio in caso di match e una penalità in caso contrario

\[
    \delta(a,b)=
    \begin{cases}
        +\mu, & a=b, \qquad a,b \in \Sigma_{\text{DNA}} \\
        -\eta, & a\neq b, \qquad a,b \in \Sigma_{\text{DNA}}\\
        -\sigma, & \text{se esattamente uno tra }a, b \text{ è un gap}
    \end{cases}
\]
con $\mu,\eta,\sigma > 0$.
Il modello in cui mismatch e gap hanno la stessa penalità si ottiene come caso particolare ponendo $\eta = \sigma$.

Grazie alla funzione di scoring possiamo anche definire la \textbf{scoring matrix}, 
\[
    \begin{array}{c|ccccc}
    \delta & A & C & G & T & -\\
    \hline
    A & +\mu & -\eta & -\eta & -\eta & -\sigma\\
    C & -\eta & +\mu & -\eta & -\eta & -\sigma\\
    G & -\eta & -\eta & +\mu & -\eta & -\sigma\\
    T & -\eta & -\eta & -\eta & +\mu & -\sigma\\
    - & -\sigma & -\sigma & -\sigma & -\sigma & \text{n.a.}
    \end{array}
\]

Siano $x,y$ le sequenze allineate corrispondenti rispettivamente a $v$ e $w$,\\
quindi $x_{i},y_{i}\in \Sigma_{\text{DNA}} \cup \{ - \}$, $\forall i \in \{ 1,\ldots,L \}$.

Sia
\[
    \mathcal{A} = \begin{pmatrix}
        x_{1} & \cdots & x_{L} \\
        y_{1} & \cdots & y_{L}
    \end{pmatrix},
\]
allora, definiamo lo score complessivo dell'allineamento come
\[
    \text{score}(\mathcal{A}) = \sum^L_{i=1} \delta(x_i,y_i).
\]

Nel modello precedente ogni gap viene penalizzato indipendentemente.
Una successione di gap consecutivi può tuttavia essere interpretata come risultato di un unico evento di inserzione o delezione.
Per questo motivo vengono spesso utilizzate penalità affini, che distinguono il costo di apertura di un gap da quelo della sua estensione.
Tale generalizzazione richiede stati aggiuntivi nella programmazione dinamica; nel seguito ci concentreremo tuttavia sul modello a penalità
lineare utilizzato dagli algoritmi oggetto del confronto.

\subsection{Programmazione dinamica}

La \emph{programmazione dinamica} è una tecnica algoritmica adatta a problemi caratterizzati da sottoproblemi sovrapposti
e da una proprietà di sottostruttura ottima.
I risultati dei sottoproblemi vengono memorizzati, evitando di calcolarli più volte.

Nel problema di allineamento che considereremo, i sottoproblemi possono essere indicizzati mediante due indici $i$ e $j$.
In particolare, una ricorrenza avrà tipicamente forma
\[
    P(i,j) =
    F\bigl(P(i-1,j),P(i,j-1),P(i-1,j-1)\bigr).
\]

Siano $|v| = n$ e $|w| = m$. Indicheremo con
\[
    S_{i,j}
\]
il valore ottimo relativo ai prefissi
\[
     v_{1}\cdots v_{i} \qquad \text{e} w_{1} \cdots w_{j},
\]
con $0 \leq i \leq n$ e $0 \leq j \leq m$

Nel caso dell'allineamento di sequenze, lo stato $S_{i,j}$ dipenderà dai tre stati
\[
    S_{i-1,j},\qquad S_{i,j-1},\qquad S_{i-1,j-1},
\]
che corrisponderanno rispettivamente a una delezione, un'inserzione e a un match o mismatch.

Memorizzando esplicitamente tutti gli stati
è necessaria una matrice di $(n+1)(m+1)$ elementi, e quindi uno spazio $O(nm)$

Una volta calcolato il valore ottimo, una soluzione corrispondente può essere ricostruita memorizzando le scelte effettuate
e percorrendole a ritroso mediante \emph{traceback} 

\subsection{Longest Common Subsequence}

Siano $v:=v_{1}\cdots v_{n}$ e $w:=w_{1}\cdots w_{n}$ due stringhe.
Una \emph{common subsequence} di $v$ e $w$ è una sequenza di lunghezza $k$, identificata dagli indici
\[
    \begin{aligned}
        1 \leq i_{1} < i_{2} < \cdots < i_{k} \leq n \\ 
        1 \leq j_{1} < j_{2} < \cdots < j_{k} \leq m
    \end{aligned}
\]
tali che
\[
    v_{i_{t}}=w_{j_{t}}\qquad \forall t \in \{ 1,\ldots,k \}
\]


Introduciamo allora tre stati per modellare la situazione appena trattata:
\begin{itemize}
    \item[(1)] $S_{i,j}$: miglior allineamento dei prefissi che termina normalmente;
    \item[(2)] $S_{i,j}^{\downarrow}$: miglior allineamento che termina con una delezione, quindi con un gap nella seconda sequenza;
    \item[(3)] $S_{i,j}^{\rightarrow}$ : miglior allineamento che termina con una inserzione, quindi con un gap nella prima sequenza.
\end{itemize}

Le transizioni dei due stati di gap, $(1)$ e $(2)$, distinguono apertura e continuazione:
\[
    S^{\downarrow}_{i,j}=\max\left\{
    S^{\downarrow}_{i-1,j}-\sigma,
    S_{i-1,j}-(\rho+\sigma)
    \right\},
\]
\[
    S^{\rightarrow}_{i,j}=\max\left\{
    S^{\rightarrow}_{i,j-1}-\sigma,
    S_{i,j-1}-(\rho+\sigma)
    \right\}.
\]

Lo stato principale $(3)$ combina quindi il match/mismatch con i due stati di gap:
\[
    S_{i,j}=\max\left\{
    S_{i-1,j-1}+\delta(v_i,w_j),
    S^{\downarrow}_{i,j},
    S^{\rightarrow}_{i,j}
    \right\}.
\]

Avendo definito il sistema di valutazione e di allineamento, possiamo studiarne la complessità


\subsection{Alignment graph e Manhattan Tourist Problem}

\subsection{Global Alignment Problem}

Delineamo gli obiettivi computazionali che ci guideranno nel percorso:
\begin{itemize}
    \item calcolare il valore ottimo;
    \item ricostruire una soluzione ottima;
    \item riuscire a farlo utilizzando poca memoria.
\end{itemize}


